% Moved from Section 6 of the long version (sections/s6_basinaware.tex) to keep the article at standard length.
% Text copied verbatim; only references are adapted (M- = main article; no prefix = Supplementary Material).

\subsection{Formulation: full text}\label{supp:std-ba-form}

The schedule of Section~\ref{M-sec:feasibility} takes the largest gain whose steady command stays within the clip at every depth up to $\kappa\alim$; the residual authority that it leaves, and with it the disturbance that the cut tolerates, is whatever results. Sections~\ref{M-sec:bistability} and~\ref{M-sec:certificates} showed that this residual authority sets the basin. Section~\ref{M-sec:basin-aware} turns the tolerated disturbance into a design requirement.

At a spindle speed $n$, a gain $g$ and a depth $a$ are called admissible if
\begin{equation*}
  a<\alim(n,g),\quad \bar b(n,a,g)\le(1-\rho)\Vmax,\quad \beta(n,a,g)\ge\beta_{\mathrm{req}},\quad a\le a_{\mathrm{cap}},
  \tag{\ref{M-eq:admissible}}
\end{equation*}
that is, if the cut is linearly stable, its forced command keeps a voltage margin $\rho$, the certified kick $\beta$ of the averaged collocation model (C2, Section~\ref{M-sec:cert-c2}) is at least the required disturbance tolerance $\beta_{\mathrm{req}}$, and the depth does not exceed an assumed cap $a_{\mathrm{cap}}$ on the axial depth. The margin leaves room for runout, noise and errors of the model in the forced command (Section~\ref{M-sec:delayed-robust}); the certificate, however, is computed for the nominal forced command with the full clip, and it shrinks if these errors consume the margin. The basin-aware depth is
\begin{equation*}
  a_{\mathrm{rob}}(n;\beta_{\mathrm{req}})=\max_{g\in\mathcal G}\max\bigl\{a\in\mathcal A:\text{ all }a'\in\mathcal A\text{ with }a'\le a\text{ admissible with }g\bigr\},
  \tag{\ref{M-eq:arob}}
\end{equation*}
with $\mathcal A$ a grid of depths and $\mathcal G$ the gains evaluated by the search (a grid refined around the best gain, Section~\ref{M-sec:ba-res}), so that $a_{\mathrm{rob}}$ and its gain are the best among these gains, not a continuous optimum, and the gain that attains it is the basin-aware gain $g_{\mathrm{rob}}(n;\beta_{\mathrm{req}})$; where several evaluated gains attain it, the one with the larger smallest certified kick up to $a_{\mathrm{rob}}$, and then the smaller gain, is taken. As neither the forced command nor the certificate is monotone in the depth, every shallower depth of the grid is required, so that, as below a stability lobe, every grid depth up to $a_{\mathrm{rob}}$ is admissible at that speed with one gain, for the settled cut (the entry is treated in Section~\ref{M-sec:ba-valid}); the depths between grid points are linearly stable and, checked every \SI{0.25}{\milli\metre}, keep the voltage margin, but are not certified. Equation~\eqref{M-eq:feas} chooses the gain for the largest linear limit under a constraint on the steady command; Eq.~\eqref{M-eq:arob} chooses it for the largest depth at which a disturbance of a given size is certified to be absorbed. Here $\rho=0.1$ and $a_{\mathrm{cap}}=2D=\SI{32}{\milli\metre}$. This cap is an assumption: the case study gives neither the cutting length of the tool nor the mode shape over such depths, and its tests went up to \SI{10}{\milli\metre}; Section~\ref{M-sec:ba-compare} repeats the chart with caps of 10 and \SI{16}{\milli\metre}.

The requirement $\beta_{\mathrm{req}}$ is an instantaneous velocity perturbation at any phase of the tooth period. As an equivalent displacement $\beta_{\mathrm{req}}/\omega_n$, 1, 2 and \SI{5}{\milli\metre\per\second} are 1.2, 2.5 and \SI{6.2}{\micro\metre}, and as an impulse $m\beta_{\mathrm{req}}$ at the four speeds of Section~\ref{M-sec:bistab-basin} they are 0.5--1.3, 1.0--2.5 and 2.4--\SI{6.3}{\percent} of the cutting impulse of one tooth period at \SI{16}{\milli\metre} [Eq.~\eqref{M-eq:impulse}]; near the linear limit the simulated basin of the DVF cut at $g=1500$ is 2.3--\SI{10.9}{\milli\metre\per\second} (Table~\ref{tab:basin}).
% src: beta / omega_n with omega_n = 2 pi x 129.3 rad/s: 1.231, 2.462, 6.154 um; impulse shares m beta / (F_y T) from Eq. (impulse) with src/params.py at 16 mm, 1090/1420/2500/2870 rpm: 0.48/0.62/1.09/1.25 %, 0.95/1.24/2.18/2.51 %, 2.38/3.10/5.46/6.27 % (recomputed); results/c1_basin.npz (worst-phase dv at 0.98 a_lim 2.34-10.88 mm/s)

\subsection{The basin-aware stability chart: full text}\label{supp:std-ba-chart}

Figure~\ref{M-fig:ba} shows the chart at 21 speeds from 500 to \SI{3000}{\rpm}, for gains of 0--\SI{3000}{\volt\second\per\metre}, searched on a grid of \SI{250}{\volt\second\per\metre} refined by $\pm125$ and $\pm\SI{62.5}{\volt\second\per\metre}$ around the best gain, and depths on a \SI{2}{\milli\metre} grid; Table~\ref{M-tab:ba} compares it with the reference schedules (Section~\ref{M-sec:ba-compare}). The certificates were scanned with a reduced collocation order and every design point, every lower depth with a certified kick below $1.15\,\beta_{\mathrm{req}}$ and the failing depth were certified again with the default order; no design point had to be lowered (Sections~\ref{supp:ba-protocol} and~\ref{supp:ba-table}).
% src: verification/paper/basin_aware.py/.json (protocol: GRID_RPM 500:125:3000, G_COARSE 0:250:3000 with refinement +-125, +-62.5, DEPTHS 2:2:32 mm, rho 0.1, a_phys 32 mm, BETA_REQ 1, 2, 5 mm/s; summary.*.corrected_down_speeds [] and upward_possible_speeds []; nred_vs_default: reduced/default 0.994-1.044)

Three features stand out. First, the linear limit at the chosen gain, not the voltage, binds at most speeds: the next grid depth is linearly unstable at 15, 17 and 15 of the 21 speeds for $\beta_{\mathrm{req}}=1$, 2 and \SI{5}{\milli\metre\per\second}, the depth cap binds at low speeds, where gains of 2500--\SI{3000}{\volt\second\per\metre} are usable, and the certificate binds at six speeds for \SI{5}{\milli\metre\per\second}, at four of which a gain skipped by the search reaches the same depth with a larger kick and is limited by the linear limit instead (Section~\ref{supp:ba-cap}). The voltage margin never binds directly, because the gain is chosen to respect it, but it is what stops a larger gain at two to four speeds. Second, the basin-aware gain is never above the gain of the authority rule with the same range and voltage margin (Fig.~\ref{M-fig:ba}c): it equals it where both reach \SI{3000}{\volt\second\per\metre}, at 500, 625 and \SI{875}{\rpm} for \SI{1}{\milli\metre\per\second} and at 500 and \SI{625}{\rpm} for \SI{2}{\milli\metre\per\second}, and is lower elsewhere; over the 21 speeds the median ratio is 0.94, 0.85 and 0.72 for the three requirements; the authority rule takes the largest gain whose steady command fits the voltage, the chart stops where the residual authority still certifies $\beta_{\mathrm{req}}$. Against the former schedule of Section~\ref{M-sec:feasibility} ($g\le\SI{1500}{\volt\second\per\metre}$, no margin) the basin-aware gain is lower between 1500 and \SI{2500}{\rpm}, where that schedule brings the steady command to the full \SI{9}{\volt} and, held to the chart's margin and certificate, keeps them only up to \SI{4}{\milli\metre}, as its forced command passes \SI{8.1}{\volt} near \SI{6}{\milli\metre}; at low speeds, and for 1 and \SI{2}{\milli\metre\per\second} also at high speeds, it is higher, because a larger gain raises the linear limit faster than it consumes the residual authority. Third, the chart keeps no explicit margin below the linear limit: the median ratio of $a_{\mathrm{rob}}$ to $\alim(g_{\mathrm{rob}})$ is 0.97, 0.97 and 0.91 for the three requirements, whereas the checked depth of Section~\ref{M-sec:feasibility} stays below $\kappa=0.9$ times its linear limit. Section~\ref{M-sec:ba-compare} adds that margin to the chart and compares it with the reference schedules.
% src: verification/paper/basin_aware.json (summary.1/2/5: primary_binding_counts_grid21 {linear 15, a_phys 6}, {linear 17, a_phys 3, certificate 1}, {linear 15, certificate 6}; gain_limited_by_primary_grid21 voltage margin 4/3/2; g_star_by_rpm; per_speed.*.auth_under_basin_criteria a_rob 4 mm at 1500-2500 rpm; per_speed.<1500..2500>.evaluated[<auth_g1500.g>].beta_1: a_rob_mm 4, fail_depth_mm 6, fail_bbar_V 8.57-9.00; fair_designs.json speeds.*.auth10.g against g_star_by_rpm: never below it, equal at 500/625/875 (1 mm/s) and 500/625 (2 mm/s), median ratio g_rob/g_auth10 0.941/0.846/0.716 (recomputed); ba_cap.json reproduction_cap32.differences (5 mm/s at 500, 625, 1000, 1375 rpm: same a_rob, binding linear limit, gain skipped by basin_aware); a_rob_mm_grid21 mean 21.71/19.90/16.00; a_chk_mm_grid21_mean 15.50; auth_g3000 a_chk mean 17.13 (recomputed over the 21 speeds); a_rob_below_a_chk_speeds 6/7/11; per_speed.*.designs.*.a_rob_over_a_lim_at_g median 0.973/0.970/0.909 over the 21 speeds, recomputed; designs.*.voltage_margin_violated_between_grid_depths false at all 63 designs; g_star_by_rpm against auth_g1500.g: 5 mm/s g_rob <= g_auth from 1000 rpm)

\paragraph{Figure~\ref{M-fig:ba} in detail} Basin-aware stability chart at the 21 speeds (markers; lines guide the eye; only the depths of the \SI{2}{\milli\metre} grid are certified). (a)~Largest depth $a_{\mathrm{rob}}$ up to which every grid depth is linearly stable, keeps \SI{10}{\percent} of the voltage free and has a certified kick (C2, averaged collocation model) of at least $\beta_{\mathrm{req}}=1$, 2 or \SI{5}{\milli\metre\per\second}, with the condition that fails at the next grid depth (markers); crosses: the largest grid depth up to which, in addition, the step entry from rest settles in the simulator with the same gain (verified numerically, not certified); the uncontrolled limit $\alim(0)$, the depth $a_{\mathrm{circ}}$ below which the averaged delay model attracts every solution at the gains of Section~\ref{M-sec:feasibility} (Proposition~\ref{M-prop:circle}; not computed at the chart's gains) and the assumed depth cap $a_{\mathrm{cap}}$. (b)~With the same gain range ($g\le\SI{3000}{\volt\second\per\metre}$) and depth margin ($a\le0.9\,\alim$): the basin-aware depth (circles), its step-entry depth (crosses) and the step-entry-checked depths $\achk$ of the authority rule with the same \SI{10}{\percent} voltage margin, Eq.~\eqref{M-eq:feas}, and of the published condition, Eq.~\eqref{M-eq:pub}, which has none. (c)~Gains of the basin-aware chart of (a), of the authority rule with the \SI{10}{\percent} margin and of the published condition.

\subsection{Comparison with the reference schedules: full text}\label{supp:std-ba-compare}

The reference schedule of Section~\ref{M-sec:feasibility} was limited to $g\le\SI{1500}{\volt\second\per\metre}$, the chart to \SI{3000}{\volt\second\per\metre}. Table~\ref{M-tab:ba} therefore compares the chart with reference schedules computed at the same 21 speeds, with the same gain range and the same rule for the operating depth, the checked depth $\achk=\min\{0.9\,\alim,a_{\mathrm{TD}}\}$ of Eq.~\eqref{M-eq:achk} at each design's own gain: the published condition, Eq.~\eqref{M-eq:pub}, reproduced as in Section~\ref{M-sec:feas-pub}, and the authority rule, Eq.~\eqref{M-eq:feas}, without voltage margin and with the \SI{10}{\percent} margin of the chart. With $g\le\SI{3000}{\volt\second\per\metre}$ the authority rule reaches 17.1 and \SI{16.8}{\milli\metre} on average instead of \SI{15.5}{\milli\metre}, so part of the gap between the chart and the former reference was the gain range. Judged by the chart's own conditions, the reference gains certify less, by an amount that depends on the requirement. The gain of the authority rule with the same voltage margin keeps every grid depth admissible up to 21.2, 18.2 and \SI{9.8}{\milli\metre} on average for 1, 2 and \SI{5}{\milli\metre\per\second}, against 21.7, 19.9 and \SI{16.0}{\milli\metre} for the basin-aware gain. These depths count the voltage margin as kept where the rigorous bound of the authority gain exceeds \SI{8.1}{\volt} by at most \SI{2.3}{\milli\volt}, at one grid depth at five speeds: the rule sets its steady command to \SI{8.1}{\volt} on sampled depths, so that this is a numerical tie, which no gain of the chart meets; counted strictly, the authority gain stops at 2 or \SI{4}{\milli\metre} there and certifies 17.9, 15.5 and \SI{9.2}{\milli\metre} on average. The other reference gains reach at most \SI{11}{\milli\metre}, because their forced command uses the whole voltage (no margin) or their gain is small. With the depth margin $a\le0.9\,\alim$ imposed on the chart as well, which then binds at 10 to 13 speeds, the chart's mean depth becomes 19.7, 18.8 and \SI{15.1}{\milli\metre}, against 19.6, 17.1 and \SI{9.4}{\milli\metre} for the authority gain under the same conditions. The choice of the gain therefore matters little for a requirement of \SI{1}{\milli\metre\per\second}, where the authority gain certifies as deep, and much for \SI{5}{\milli\metre\per\second}, where it certifies \SI{5.7}{\milli\metre} less: the authority rule spends the voltage on raising the linear limit, the chart keeps the residual authority that a larger disturbance needs. The chart is a grid optimum, however: at five speeds between 1625 and \SI{2375}{\rpm} the authority gain, which lies between the gains the chart evaluated, meets the chart's conditions \SI{2}{\milli\metre} deeper for some of the requirements, and with the depth margin at 1375, 1625 and \SI{2625}{\rpm}. Against the authority rule's own operating depth with the same margins (\SI{16.8}{\milli\metre}), the chart is 2.9 and \SI{2.0}{\milli\metre} deeper for 1 and \SI{2}{\milli\metre\per\second} and \SI{1.6}{\milli\metre} shallower for \SI{5}{\milli\metre\per\second}, where its advantage is the certified disturbance, not the depth.

At their own operating depths the references keep larger certified kicks than the chart for 1 and \SI{2}{\milli\metre\per\second}, because they cut less deep (Table~\ref{M-tab:ba}): the authority rule with the margin, at \SI{16.8}{\milli\metre}, has a certified kick of \SI{7.9}{\milli\metre\per\second} and a simulated critical kick of \SI{15.1}{\milli\metre\per\second} (medians over the speeds), against 4.3 and \SI{8.1}{\milli\metre\per\second} for the chart for \SI{1}{\milli\metre\per\second} at \SI{21.7}{\milli\metre}; with the depth margin, at \SI{19.7}{\milli\metre}, the chart's certified kick stays at \SI{4.3}{\milli\metre\per\second} but its critical kick, \SI{16.4}{\milli\metre\per\second}, is larger. For \SI{5}{\milli\metre\per\second}, at nearly the same depth, the chart's kicks are similar, 7.2 and \SI{15.4}{\milli\metre\per\second}. What the authority rule does not check is the shallower depths that its gain passes through: the smallest certified kick of its gain over the grid depths below its operating depth and at it is \SI{2.3}{\milli\metre\per\second} (median) and reaches 1, 2 and \SI{5}{\milli\metre\per\second} at 20, 18 and 5 of the 21 speeds, whereas the chart certifies $\beta_{\mathrm{req}}$, on the averaged collocation model, at every grid depth up to $a_{\mathrm{rob}}$. The chart thus moves the operating point to the deepest grid depth at which the required kick is still certified, while the reference keeps a reserve that it does not quantify. The published condition, at the smallest depths, keeps the largest kicks, but its forced command exceeds the clip at \SI{2000}{\rpm} ($\bar b=\SI{9.9}{\volt}$), and without voltage margin the authority rule and the former reference reach the clip at one speed each, where no certificate exists. The step entries of the references settle at their operating depths at all 21 speeds, those of the chart at 7, 9 and 14 speeds (Section~\ref{supp:std-ba-valid}).
% src: verification/paper/fair_designs.json (summary); verification/paper/ref_under_ba.json, ref_under_ba_kappa0.9.json (summary); verification/paper/ba_kappa.json (summary: binding_counts depth margin 13/13/10); ref_under_ba_tol.json speeds.*.auth10.tol_k1/tol_k0.9 against basin_aware.json and ba_kappa.json a_rob (2 mm above the chart at 1625 rpm for 1 mm/s, 1875 and 2000 rpm for 1, 2, 5 mm/s, 2125 and 2375 rpm for 1, 2 mm/s; with the depth margin at 1375 (1), 1625 (1, 2) and 2625 rpm (1, 2); recomputed); ba_kappa.json summary.mean_a_rob_mm 19.71/18.76/15.14 against fair_designs.json mean_a_chk_mm.auth10 16.77: +2.94/+1.99/-1.63 mm
% src (kicks at the operating depth): verification/paper/fair_compare.py/.json (summary.auth10: beta_op median 7.92, crit_kick median 15.06, beta_range median 2.29 (minimum over the grid depths below a_op and a_op itself), beta_range_at_least 1/2/5: 20/18/5 of 21; table6.json rows_kappa.ba1: beta median 4.27, crit median 16.42, depth 19.71; ba1: beta_op 4.27, crit 8.12, a_op mean 21.71; ba5: beta_op 7.20, crit 15.39; pub: bbar max 9.90 at 2000 rpm; auth0 and auth1500: bbar max 9.00 (at 1625 rpm), beta_op min 0; step_entry_stable: references 21, ba1/ba2/ba5 7/9/14; checks.recheck: stored kick and certificate recomputed bit for bit); TIGHT over FAST at most 1.004 over the 17 points with both (recomputed), so the FAST-only chart values (checks.ba_beta_op_fast_only_below_5) are not disadvantaged

The assumed cap $a_{\mathrm{cap}}=\SI{32}{\milli\metre}$ binds at 6, 3 and 0 of the 21 speeds for the three requirements, but the deep cuts it admits carry much of the depth advantage. With the chart repeated for caps of 16 and \SI{10}{\milli\metre}, the latter the deepest of the published tests (Section~\ref{supp:ba-cap}; with \SI{32}{\milli\metre} the repetition returns the stored depths at 62 of the 63 design points), the chart reaches \SI{16}{\milli\metre} at 16, 16 and 9 speeds, with mean depths of 15.1, 15.1 and \SI{14.0}{\milli\metre}, and \SI{10}{\milli\metre} at every speed. Capped the same way, at $\min\{\achk,a_{\mathrm{cap}}\}$, the authority rule with the same voltage margin operates at 14.5 and \SI{10.0}{\milli\metre} on average and the published condition at 10.7 and \SI{8.9}{\milli\metre}. Below a cap of about \SI{16}{\milli\metre} the chart therefore gains little or no depth over the authority rule: \SI{0.6}{\milli\metre} for 1 and \SI{2}{\milli\metre\per\second} and \SI{0.5}{\milli\metre} less for \SI{5}{\milli\metre\per\second} with \SI{16}{\milli\metre}, none with \SI{10}{\milli\metre}. Over the authority gain it still gains certified disturbance at equal depth. At the speeds where both gains meet the chart's conditions up to the cap, the smallest certified kick over the grid depths is larger with the chart's gain than with the authority gain by a median factor of 2.3 for \SI{10}{\milli\metre} (21 speeds) and 1.5 for \SI{16}{\milli\metre} (16 speeds), with the same tolerance on the voltage test. Against the former reference the factor is 1.24 and 1.00, and against the published gain, where it reaches \SI{10}{\milli\metre} (12 speeds), 1.02; ratios within about \SI{12}{\percent} of one are ties, because the reference gains are judged with the reduced collocation order, which overstates the default-order kick by up to \SI{12}{\percent} there. With a shallow cap the chart also chooses low gains (\SI{700}{\volt\second\per\metre} on average for \SI{10}{\milli\metre}).
% src: results/cuts_main.json (deepest published cut 10 mm); verification/paper/ba_cap.py/.json (summary.per_cap.32.*.n_at_cap 6/3/0; per_cap.16.*: mean_a_rob_mm 15.14/15.14/14.00, n_at_cap 16/16/9; per_cap.10.*: mean_a_rob_mm 10.00, n_at_cap 21, mean_g 699.4; per_cap.<cap>.ref_capped_mean_mm: 16 -> auth10 14.54, pub 10.69; 10 -> auth10 10.00, pub 8.88; reproduction_cap32: same_a_rob 62 of 63); verification/paper/cap_kicks.py/.json and cap_kicks_tol.json (--tol: reference scans of ref_under_ba_tol.py; caps.<cap>.paired_both_at_cap: 10 auth10_vs_ba_1 n 21 median 2.30 (strict n 16, 2.09), auth1500 n 12 1.24, pub n 12 1.02; 16 auth10_vs_ba_1 n 16 median 1.52 (strict n 13, 1.38), auth1500 n 8 1.00; ba_cap.json summary.per_cap.16.<b>.a_rob_minus_ref_capped_mean_mm.auth10 +0.60/+0.60/-0.55, per_cap.10 0.0; reduced over default order at the reference operating depths 0.975-1.122, median 1.008 (fair_compare.json per_speed beta_range at a_op against beta_op.fast_mmps, recomputed))

\subsection{Validation, entry and robustness: full text}\label{supp:std-ba-valid}

At six speeds (750, 1125, 1375, 2000, 2500 and \SI{2875}{\rpm}) the design points were simulated with the \SI{9}{\volt} clip and the full surface memory. The simulated critical kick, smallest over four phases a quarter of a tooth period apart, exceeds $\beta_{\mathrm{req}}$ by factors of at least 2.75, 2.78 and 2.27 for 1, 2 and \SI{5}{\milli\metre\per\second}, or 1.6--6.9 times the certified kick, and kicks of exactly $\beta_{\mathrm{req}}$ at all 16 phases never caused chatter (0 of 96 for each requirement). Computed at all 21 speeds (Table~\ref{M-tab:ba}), the smallest critical kick still exceeds $\beta_{\mathrm{req}}$ by factors of at least 2.55, 1.97 and 2.08, and with the depth margin 4.20, 3.49 and 2.30: in the simulator no design point of the chart has a smaller basin than required, along the kick direction.
% src: verification/paper/basin_aware.json (validation_summary.1/2/5: points 6, all_above_beta_req true, worst_crit_kick_over_beta_req_min 2.750/2.780/2.272, td_over_cert 1.63-6.95, kicks_of_beta_req_chattering 0 of 96); verification/paper/fair_compare.json (summary.ba1/ba2/ba5.crit_kick_mmps.min 2.549/3.931/10.419, i.e. 2.55/1.97/2.08 times beta_req; crit_kick_unsettled 0) and bak_crit.json (min 4.195/6.979/11.486, i.e. 4.20/3.49/2.30)

The chart certifies the settled cut, not the entry. With the gains of the chart, the face entry settles at every grid depth up to $a_{\mathrm{rob}}$ at all 21 speeds, but the step entry from rest only up to 16.1, 15.8 and \SI{15.1}{\milli\metre} on average for the three requirements (crosses in Fig.~\ref{M-fig:ba}a), below $a_{\mathrm{rob}}$ at 14, 12 and 7 speeds, and with the depth margin up to 15.7, 15.4 and \SI{14.4}{\milli\metre}. These depths are verified in the simulator, not certified, and a few step entries counted as settled pass through transients of some millimetres; excluding those above \SI{1}{\milli\metre} lowers the averages by at most \SI{0.4}{\milli\metre}. With an abrupt entry the usable depth is thus set by saturation during the entry, for the chart as for the authority rule, whose checked depth includes the same check (\SI{16.8}{\milli\metre} with the same margins, Fig.~\ref{M-fig:ba}b); a certificate of the entry itself, along $-\xi_p(0)$, computed at $g=\SI{1500}{\volt\second\per\metre}$ at the four speeds of Section~\ref{M-sec:bistab-basin}, succeeded only at $0.3\,\alim$ (Section~\ref{M-sec:cert-c2}) and was not computed at the chart's design points.
% src: verification/paper/ba_entry.py/.json (summary.*: n_a_face_below_a_rob 0; mean_a_step_mm 16.10/15.81/15.14; n_a_step_below_a_rob 14/12/7; quiet.mean_a_step_quiet_mm 15.71/15.52/15.14), ba_entry_kappa.json (mean_a_step_mm 15.71/15.43/14.38; quiet 15.33/15.14/14.38); verification/paper/fair_designs.json (summary.mean_a_chk_mm auth10 16.77)

Under errors in the model the operating points of the chart, closer to their own linear limit, are less robust than the reference. The first 30 samples of the $\pm\SI{10}{\percent}$ error draw of Section~\ref{M-sec:delayed-robust}, the same at every speed and for every design, were applied at the six validation speeds to the chart and to the authority rule with the \SI{10}{\percent} margin (Section~\ref{supp:ba-mc}). A sample fails if the cut does not settle or a kick of $\beta_{\mathrm{req}}$ at one of four phases ends in chatter; if the outcome is monotone in the kick, this is the event $\Delta v^*\le\beta_{\mathrm{req}}$, with $\Delta v^*$ the smallest critical kick over these phases. At their own depths the failure probability is 16.7, 11.7 and \SI{10.0}{\percent} for the chart with the depth margin and 39, 39 and \SI{11}{\percent} without it, for 1, 2 and \SI{5}{\milli\metre\per\second}, against \SI{6.7}{\percent} for the reference; no sample failed for the reference and not for the chart (McNemar $p\le0.031$, pooled over the speeds and approximate because they share the samples; with the margin, exact sign test over the samples $p=0.002$, 0.03 and 0.13). For \SI{5}{\milli\metre\per\second} the chart is on average shallower than the reference at these speeds and still fails more often: what matters is the ratio to the linear limit, 0.81--0.90 for the chart with the depth margin and 0.86--0.996 without it, against 0.51--0.81 for the reference where its step-entry check binds and 0.90 at \SI{2000}{\rpm}, where its depth margin binds and where alone it fails, in the same samples as the chart. Of the 69 failures of the chart with the depth margin, 55 are cuts whose sampled averaged model is linearly unstable and 10 are cuts that do not settle although that model is stable, all at \SI{2875}{\rpm}, where the forced command of the settled cut was on the sample's clip in 18 of the 30 samples for \SI{1}{\milli\metre\per\second}; a settled cut chattered after the kick in 1, 1 and 2 of the 180 sample--speed pairs. The errors therefore consume mainly the margin to the linear limit, at \SI{2875}{\rpm} also the voltage margin, and only rarely the certified basin. At the common depth $\min\{\achk,a_{\mathrm{rob}}\}$ the probabilities are 7.2, 7.2 and \SI{9.4}{\percent} against \SI{6.1}{\percent}; the six additional failures for \SI{5}{\milli\metre\per\second} ($p=0.13$ over the samples) are linear instabilities at 1125, 2000 and \SI{2500}{\rpm}, where the chart's gains are lower, and three failures at \SI{2875}{\rpm}, two of them after the kick. These probabilities, with pooled \SI{95}{\percent} Clopper--Pearson intervals of about $\pm4$ to $\pm7$ percentage points (approximate, as the speeds share the samples), describe the sampled distribution and guarantee nothing; they indicate that the chart needs margins that cover the errors of the model, below the linear limit and, where runout and a lower clip consume the margin of the forced command, in the voltage, which the certificate of the nominal model does not provide.
% src: verification/paper/mc_designs.py (protocol), mc_designs_kappa.json and mc_designs_chart.json (complete; summary.<b>.own/common: ref 12/180 (own), 11/180 (common), cp95 [0.035, 0.114]; chart with margin own 30/21/18 of 180 (0.167/0.117/0.100; unsettled 29/20/16, kick 1/1/2; McNemar 7.6e-6/0.0039/0.031; sign 0.0020/0.031/0.125); common 13/13/17 (0.072/0.072/0.094; sign 0.5/0.5/0.125); chart without margin own 71/71/19 (0.394/0.394/0.106; McNemar 3.5e-18/3.5e-18/0.016); ref_fails_BA_passes 0 everywhere); verification/paper/mc_designs_table.py/.json (rows: a_over_a_lim ref 0.512-0.810 at the five speeds where a_td < 0.9 a_lim and 0.900 at 2000 rpm (fair_designs.json speeds.2000.auth10: a_td = kappa_a_lim = a_chk 15.25 mm), kappa 0.809-0.899, chart 0.862-0.996; 5 mm/s chart with margin mean depth 14.67 mm against reference 15.90 mm at the six speeds; 2000 rpm ref g 213.7, 12/30 failures, both_fail 12); mc_designs_kappa.json per_speed.<rpm>.b.<b>.own.BA: sums of unsettled_and_cf_unstable 21/20/14 (55), failed_but_cf_stable 9 (2875: 8 unsettled, 1 kick)/1 (750: kick)/4 (2875: 2 unsettled, 2 kick), settled_pre_kick_clipped 18 at 2875 rpm for 1 mm/s; common b 5: BA_fails_ref_passes 6 (1125, 2000, 2500 rpm linearly unstable; 2875 rpm 3, kick 2)
